\documentstyle[%


righttag%


,11pt%


,amssymb%


,russian%               remove in Eng.


,izvru%                 Set izven% in Eng.


]{amsart}





% 





\renewcommand{\baselinestretch}{1.02}


\newcommand{\re}{\operatorname{Re}}


\newcommand{\im}{\operatorname{Im}}


\newcommand{\tts}[1]{}





\theoremstyle{definition}


\theorembodyfont{\rm}


\newtheorem{defnnonum}{\hskip\parindent Определение}


\newtheorem{cornonum}{\hskip\parindent Следствие}


\renewcommand{\thecornonum}{}


\renewcommand{\thedefnnonum}{}


\renewcommand{\baselinestretch}{1.03}





%\hoffset=-15mm%                  For Eng.


\setcounter{page}{28}


\god{2004}


\nomer{4~(503)} %                        Remove in Eng.


%\nomer{Vol.\,48, No.\,4}%               Set in Eng.


%\str{\pageref{begin}--\pageref{end}}%  Set in Eng.


\udk{517.546}





\author{\it В.В.\,Кожевников}


% NB:   ^^ remove in Eng.


\title{Об одной форме прямого вариационного метода\\


для однолистных функций}





\date{}


%\pagestyle{myheadings}%         Set in Eng.


\pagestyle{plain} %              Remove in Eng.


\begin{document}


%\thispagestyle{plain}%          Set in Eng.


\maketitle


\label{begin}











В работе предлагается оригинальный подход к построению


вариационного метода для однолистных функций,


основанный на теореме В.Д.\,Ерохина о факторизации


[1], согласно которой любое конформное отображение 


конечносвязной области можно представить в виде


композиции конформных отображений односвязных областей.


Первые попытки применить теорему о факторизации 


в вариационных методах принадлежат 


С.А.\,Гельферу [2]. Используя метод Г.М.\,Голузина [3], 


он получил однолистные 


вариации для многосвязной области путем редукции 


этой задачи к односвязной ситуации. 





В данной работе однолистные вариации 


для произвольной односвязной области строятся


по однолистным вариациям, определенным в 


окрестности ее границы.


В результате удается построить весьма общую форму


однолистных вариаций. Эта форма представляет собой


только оболочку, в которую можно заключить конкретный


класс вспомогательных однолистных вариаций, параметризованный


тем или иным способом. Различные способы параметризации


приводят к различным вариационным схемам. В частности,


применение метода может строиться по схеме метода


граничных или внутренних вариаций М.\,Шиффера [4], [5],


а также вариационного метода Г.М.\,Голузина.





Особо следует упомянуть работу [6],


в которой предложена модификация метода Г.М.\,Го\-лу\-зи\-на и 


получены однолистные вариации, по форме тождественные 


вариациям (8) из данной работы, но способ получения


вариаций и характер их использования существенно иные. 





Приводимые в работе результаты ранее были опубликованы в [7], но 


без объяснений и практически без доказательств.


\medskip





{\bf 1.}


Исходным является





\begin{defnnonum}


Однозначную функцию $v=v(w,\tau)$ комплексной


переменной $w\in G$ и вещественной переменной


$\tau$, $0<\tau<\tau^{*}$, условимся называть


{\it однолистной вариацией}, определенной в области $G$, если


\begin{itemize}


\item[---] при каждом значении $\tau$ функция $v(w,\tau)$  


как функция переменной $w$ однолистна в области $G$; 


\item[---] на любом замкнутом подмножестве области


$G$ для функции $v=v(w,\tau)$ выполняется


равномерная оценка 


$$


v(w,\tau)=w+o(1),\quad \tau\to 0+;


$$


здесь $o(1)$ --- функция переменных $w$, $\tau$,


которая на любом замкнутом подмножестве области


$G$ равномерно по переменной $w$ стремится к 


нулю при $\tau\to 0+$.


\end{itemize}


Если, кроме того, вариация $v(w,\tau)$ при каждом


значении $\tau$ как функция переменной


$w$ регулярна в области $G$, то будем называть ее однолистной 


и регулярной.


\end{defnnonum}








\begin{thm}[{\series{m}\shape{n}\selectfont [7]}]


Пусть $\Gamma$ --- невырожденный


континуум, дополнение к которому состоит из 


односвязной области, и


$\nu=\nu(\omega,\tau)$, $0<\tau<\tau^{*}$, ---


однолистная вариация, регулярная в некоторой


проколотой окрестности $\Dot{U}(\Gamma)$


континуума $\Gamma$.\,\footnote{Под проколотой окрестностью континуума


$\Gamma$ понимаем любую, не содержащую бесконечности


двусвязную область, одна из двух компонент дополнения к которой 


до всей плоскости совпадает с $\Gamma$. Проколотую


окрестность в отличие от полной окрестности обозначаем 


точкой сверху.}





Тогда в дополнении к континууму $\Gamma$


существует однолистная вариация $v(w,\tau),$ для


которой на любом замкнутом подмножестве $E$ из


дополнения к континууму $\Gamma$ при достаточно


малых значениях $\tau$ выполняется


равномерная оценка


\begin{equation}


v(w,\tau)=w+


\frac{1}{2\pi i}\int_{\gamma}\frac{\nu-\omega}


{\omega-w}d\omega


+o(\|\nu-\omega\|_{\gamma}).\,\footnotemark


\end{equation}


\footnotetext{Под $\|\dotsc\|_{\gamma}$ подразумевается норма 


какого-либо функционального пространства. Норму


можно выбрать достаточно произвольно, но 


естественно использовать норму пространства $C(\gamma)$.}





\noindent{}Здесь $\gamma$ --- простая замкнутая 


спрямляемая кривая, лежащая в окрестности 


$\Dot{U}(\Gamma),$ разделяющая граничные компоненты


этой окрестности и отделяющая 


множество $E$ от континуума~$\Gamma;$ интегрирование


вдоль кривой $\gamma$ осуществляется в направлении, 


при котором континуум~$\Gamma$ остается слева$;$ 


функция $o(1),$ входящая в остаточный член формулы $(1)$,


регулярна по переменной $w$ в дополнении к континууму $\Gamma$


и при $\tau\to 0+$ равномерно стремится к нулю на множестве~$E$.


\end{thm}


\setcounter{footnote}{0}





Теорема 1 доказывается с помощью теоремы о факторизации.


Она позволяет строить однолистные вариации


в произвольной односвязной области по однолистным


вариациям, определенным в окрестности ее границы. 


\medskip





{\bf 2.}


При отыскании однолистных вариаций в окрестности границы


данной односвязной области применяем вспомогательные 


однолистные вариации, которые строятся в форме


конформных отображений колец. 





Пусть $Q(\zeta)$ --- функция, 


регулярная в кольце $D[\delta,1)=\{\zeta:\delta\le|\zeta|<1\}$, 


с неотрицательной реальной частью


на окружности $|\zeta|=\delta$.


Для любого малого $\varepsilon>0$ при достаточно 


малых значениях $k$, $0<k<k^{*}$, функция [3], [6]


\begin{equation}


\lambda(\zeta,k)=\zeta \exp\left (k Q(\zeta)\right )


\end{equation} 


однолистна в кольце $D[\delta,1-\varepsilon]=


\{\zeta:\delta\le|\zeta|\le1-\varepsilon\}$.


Действительно, обозначив через $M$ наибольшее значение 


модуля функции $Q(\zeta)$ в кольце $D[\delta,1-\varepsilon]$,


получим в этом кольце неравенство


\begin{equation}


\mathrm e^{-k^{*}M}\le \left |\frac{\lambda(\zeta,k)}{\zeta}\right |


\le \mathrm e^{k^{*}M}.


\end{equation}


Пусть $\zeta_{1}$, $\zeta_{2}$ --- две произвольные различные


точки кольца $D[\delta,1-\varepsilon]$. Ввиду (3) 


\begin{multline*}


|\lambda(\zeta_{2},k)-\lambda(\zeta_{1},k)|=


\bigg |(\zeta_{2}-\zeta_{1})\frac{\lambda(\zeta_{2},k)}


{\zeta_{2}}+


\zeta_{1}\bigg (\frac{\lambda(\zeta_{2},k)}{\zeta_{2}}-


\frac{\lambda(\zeta_{1},k)}{\zeta_{1}}\bigg )\bigg |\ge \\


%


\ge \mathrm e^{-k^{*}M}|\zeta_{2}-\zeta_{1}|-


\mathrm e^{k^{*}M}\left |\exp\left (kQ(\zeta_{2})-


kQ(\zeta_{1})\right )-1\right |.


\end{multline*}


Используя разложение экспоненты в ряд, легко получим


$$


\left |\exp\left (kQ(\zeta_{2})-


kQ(\zeta_{1})\right )-1\right |\le


k^{*}\left |Q(\zeta_{2})-Q(\zeta_{1})\right |\mathrm e^{2k^{*}M}.


$$


Теперь воспользуемся равенством


$$


Q(\zeta_{2})-Q(\zeta_{1})=


\int_{\zeta_{1}}^{\zeta_{2}}\frac{d}


{d\zeta}Q(\zeta)~d\zeta,


$$


в котором интегрирование осуществляется вдоль 


дуги $\overset{\Large\frown}{\zeta_{1}\zeta_{2}}$ окружности


с центральным углом, не превосходящим $\pi$, 


целиком лежащей в замкнутом круговом кольце 


$D[r,R]$, где $r=\min(|\zeta_{1}|, |\zeta_{2}|)$ и 


$R=\max(|\zeta_{1}|, |\zeta_{2}|)$. 


Обозначим через $M'$ максимальное


значение модуля производной функции $Q(\zeta)$ в кольце


$D[\delta,1-\varepsilon]$. На основании предыдущего равенства


$|Q(\zeta_{2})-Q(\zeta_{1})|\le\pi M'|\zeta_{2}-\zeta_{1}|$.


Сопоставляя полученные неравенства, приходим к оценке


\begin{equation}


|\lambda(\zeta_{2},k)-\lambda(\zeta_{1},k)|\ge 


(\mathrm e^{-k^{*}M}-k^{*}\pi M'\mathrm e^{3k^{*}M})|\zeta_{2}-\zeta_{1}|.


\end{equation}


Выберем $k^{*}>0$ столь малым, чтобы выполнялось неравенство


$k^{*}\mathrm e^{4k^{*}M}<\frac{1}{\pi M'}$,


тогда вследствие неравенства (4) при $0<k<k^{*}$ 


функция $\lambda(\zeta,k)$ однолистна в кольце


$D[\delta,1-\varepsilon]$.





В дальнейшем условие $0<k<k^{*}$ будем считать выполненным.





Перечислим некоторые свойства функции $\lambda(\zeta,k)$.





Прежде всего, поскольку реальная часть функции $Q(\zeta)$


на окружности $|\zeta|=\delta$ неотрицательна, функция


$\lambda(\zeta,k)$ конформно отображает окружность


$|\zeta|=\delta$ в некоторую замкнутую,


аналитическую кривую Жордана, целиком расположенную в замкнутой 


области $|\zeta|\ge \delta$, причем ориентация кривой


не изменяется, т.\,к. в силу регулярности на окружности


$|\zeta|=\delta$ функции


$Q(\zeta)$ приращение $\arg\lambda(\zeta,k)$ 


при полном обходе окружности $|\zeta|=\delta$ в положительном 


направлении составляет $2\pi$. Из этого, в частности, 


следует, что значения функции $\lambda(\zeta,k)$,  отвечающие 


точкам открытого кольца $D(\delta,1-\varepsilon)$, лежат


вне круга $|\zeta|\le\delta$.


Наконец, в кольце $D[\delta,1-\varepsilon]$


для функции $\lambda(\zeta,k)$ выполняется равномерная оценка


\begin{equation}


\lambda(\zeta,k)=\zeta+k\zeta Q(\zeta)+o(k),\quad k\to 0,


\end{equation}


причем для остаточного члена имеет место неравенство


\begin{equation}


|o(k)|\le \mathrm e^{kM}-1-kM.


\end{equation}





Наиболее общий способ построения функции $Q(\zeta)$


получается с использованием интеграла Вилля, но такой 


способ оказывается неоправданно громоздким. В пункте~4 


данной работы  будет показано, что для построения функции


$Q(\zeta)$ достаточно всего четырех элементарных функций,


чтобы получить информацию об экстремальном континууме,


которую может дать только вариационный метод первого


порядка. Эти функции --- функция поворота, степенная 


функция, функция Жуковского и, наконец, дробно-линейная 


функция.  


\medskip





{\bf 3.}


Вариации (1) представляют собой общую


форму\,\footnote{


С помощью формулы Коши произвольную вариацию 


$v(w,\tau)$, $0<\tau<\tau^{*}$, однолистную 


в дополнении к континууму $\Gamma$, можно с точностью до  


бесконечно малого сдвига представить в виде (1).}


однолистных вариаций, которую


можно использовать для построения однолистных вариаций


частного вида, более 


приспособленных для решения экстремальных задач. 


Для построения специального 


класса однолистных вариаций в окрестности границы данной


односвязной области используются


вспомогательные вариации (2). 





Приведем подробное доказательство теоремы из статьи [7].


\setcounter{footnote}{0}





\begin{thm}


Пусть $\Gamma$ --- невырожденный


континуум, дополнение к которому состоит из 


односвязной области, и


$\zeta(\xi)$ --- конформное отображение некоторой


проколотой окрестности $\Dot{U}(\Gamma)$


континуума $\Gamma$ на подходящее кольцо


$D(\delta,1)$, при котором граничной 


компоненте окрестности $\Dot{U}(\Gamma)$,


принадлежащей континууму $\Gamma$,


отвечает внутренняя граничная компонента кольца $D(\delta,1)$.


Если в окрестности $\Dot{U}(\Gamma)$


\begin{equation}


p(\xi)=\frac{i\zeta(\xi)}{\zeta'(\xi)}Q(\zeta(\xi)),


\end{equation}


где $Q(\zeta)$ --- функция, 


регулярная в кольце $D[\delta,1)$, 


с неотрицательной реальной частью


на окружности $|\zeta|=\delta$, то


в дополнении к континууму $\Gamma$


существует однолистная вариация $v(w,\tau)$, для


которой на любом замкнутом подмножестве $E$ из


дополнения к континууму $\Gamma$ при достаточно


малых значениях $\tau$, $0<\tau<\tau^{*}$, выполняется


равномерная оценка


\begin{equation}


v(w,\tau)=w+\tau\int_{\gamma}\frac{p(\xi)}


{\xi-w}d\xi+o(\tau).


\end{equation}


Как и в теореме $1$, $\gamma$ --- простая замкнутая


спрямляемая кривая, лежащая в окрестности 


$\Dot{U}(\Gamma)$, разделяющая граничные компоненты


этой окрестности и отделяющая 


множество $E$ от континуума $\Gamma;$ интегрирование


вдоль кривой $\gamma$ осуществляется в направлении, 


при котором континуум $\Gamma$ остается слева$;$


функция $o(1)$, входящая в остаточный член формулы $(8)$,


регулярна по переменной $w$ в дополнении к континууму $\Gamma$


и при $\tau\to 0+$ равномерно стремится к нулю на множестве $E$.


\end{thm}





\begin{pf}


Обозначим через $\xi(\zeta)$ функцию, 


обратную к функции $\zeta(\xi)$.


По $Q(\zeta)$ из (7) с помощью формулы (2)


построим функцию $\lambda(\zeta,k)$.


При достаточно малых значениях 


$k$, $0<k<k^{*}$, функция $\lambda(\zeta,k)$


однолистна в кольце $D(\delta,1-\varepsilon)$ и для нее


выполняется равномерная оценка (5).


Величину $\varepsilon$ с самого начала можно выбрать


столь малой, чтобы кривая $\gamma$ лежала в области


$\xi(D(\delta,1-\varepsilon))$.





При $0<k<k^{*}$ композиция


$\nu(\zeta)=\xi\circ\lambda(\zeta,k)$


регулярна и однолистна в кольце $D(\delta,1-\varepsilon)$.


Эту композицию можно рассматривать как результат


действия на функцию $\xi=\xi(\zeta)$


вариации $\nu(\xi,k)=\xi\circ\lambda(\zeta(\xi),k)$,


для которой на кривой $\gamma$ при $k\to 0+$ выполняется 


равномерная оценка


\begin{equation}


\nu(\xi,k)=\xi+k\zeta \xi'(\zeta)Q(\zeta)+o(k).


\end{equation}





Докажем оценку (9). Для этого внутри кольца 


$D(\delta,1-\varepsilon)$ возьмем некоторое замкнутое,


концентрическое с ним кольцо $R$, содержащее в своей


внутренности кривую $\beta=\zeta(\gamma)$. Пусть $d$ 


обозначает расстояние от кривой $\beta$ до границы  


кольца $R$. Положим 


$\rho(\zeta,k)=\lambda(\zeta,k)-\zeta$.


Вследствие (5), (6) для $\rho(\zeta,k)$ в кольце


$D(\delta,1-\varepsilon)$ выполняется оценка


\begin{equation}


|\rho(\zeta,k)|\le \mathrm e^{kM}-1.


\end{equation}


Будем считать, что $\mathrm e^{kM}-1<d$ при $0<k<k^{*}$. 


Для фиксированной точки $\zeta\in \beta$ при $0<k<k^{*}$


имеет место разложение в ряд Тейлора


\begin{equation}


\nu(\zeta)=\xi\circ\lambda(\zeta,k)=\xi(\zeta+\rho(\zeta,k))=


\sum_{n=0}^{\infty}\frac{\xi^{(n)}(\zeta)}{n!}\rho^{n}(\zeta,k).


\end{equation}


Отсюда


\begin{equation}


\nu(\zeta)=\xi(\zeta)+k\zeta\xi'(\zeta)Q(\zeta)+r(\zeta,k),


\end{equation}


где


$$


r(\zeta,k)=\xi'(\zeta) o(k)+


\sum_{n=2}^{\infty}\frac{\xi^{(n)}(\zeta)}{n!}\rho^{n}(\zeta,k),


$$


$o(k)$ --- остаточный член из формулы (5).


Обозначим через $M_{1}$ максимум модуля функции $\xi(\zeta)$


в кольце $R$. При $\zeta\in \beta$ имеем оценку,


не зависящую от $\zeta$, 


\begin{equation}


|\xi^{(n)}(\zeta)|=\bigg|\frac{n!}{2\pi i}\int_{\partial R}


\frac{\xi(\zeta')}{(\zeta'-\zeta)^{n+1}}~d\zeta'\bigg|


\le \frac{2M_{1}n!}{d^{n+1}}.


\end{equation}


Следовательно, учитывая (6) и (10),


при $0<k<k^{*}$ и $\zeta\in \beta$ будем иметь неравенство


\begin{equation}


\|r(\zeta,k)\|_{\gamma}\le


\frac{2M_{1}}{d^{2}}


\bigg(\mathrm e^{kM}-1-kM+


d\frac{(\frac{\mathrm e^{kM}-1}{d})^{2}}


{1-\frac{\mathrm e^{kM}-1}{d}}\bigg).


\end{equation}


Правая часть неравенства при $k\to 0+$ является бесконечно малой 


величиной порядка $k^{2}$, не зависящей от $\zeta$.


Таким образом, остаточный член $r(\zeta,k)$ в


формуле (12) является на кривой $\beta$ равномерно 


бесконечно малой величиной порядка $o(k)$ при $k\to 0+$.


В формуле (12) положим $\zeta=\zeta(\xi)$.


В результате получим оценку (9), равномерную на кривой $\gamma$.





Теперь уже нетрудно получить вариацию (8). 


Из формулы (9) выразим разность $\nu-\xi$


и подставим ее в подинтегральное выражение формулы (1).


В итоге, если в первом вариационном члене положить 


$\tau=k/2\pi$ и учесть, что


$\xi'(\zeta)=1/\zeta'(\xi)$, получим формулу (8)


с остаточным членом


\begin{equation}


r_{1}(w,\tau)=\frac{1}{2\pi i}\int_{\gamma}


\frac{r(\zeta,k)}{\xi-w}~d\xi + o(\|\nu-\xi\|_{\gamma}),


\end{equation}


где $o(\|\nu-\xi\|_{\gamma})$ --- остаточный член из формулы (1),


а $\nu=\nu(\xi,k)$ --- вариация (9). Поскольку


$\nu(\xi,k)-\xi=\nu(\zeta)-\xi(\zeta)$, $\xi=\xi(\zeta)$,


то на основании (11) и (13) при $\zeta\in\beta$ и $0<k<k^{*}$ 


находим


\begin{equation}


\|\nu-\xi\|_{\gamma}\le


\frac{2M_{1}}{d}\frac{\frac{\mathrm e^{kM}-1}{d}}


{1-\frac{\mathrm e^{kM}-1}{d}}.


\end{equation}


Пусть $d_{1}$ --- расстояние между заданным


замкнутым множеством $E$ и кривой $\gamma$. Ввиду


(15) при $w\in E$ и $\zeta\in \beta$ выполняется оценка


\begin{equation}


|r_{1}(w,\tau)|\le\frac{\|r(\zeta,k)\|_{\gamma}


L_{\gamma}}{2\pi d_{1}}+


|o(1)|\|\nu-\xi\|_{\gamma},


\end{equation}


где $L_{\gamma}$ --- длина кривой $\gamma$, а $o(1)$ ---


равномерно бесконечно малая при $\tau\to 0+$ и $w\in E$


--- величина, входящая в состав остаточного члена из формулы (1).


Из оценок (14), (16), (17) следует, что остаточный член


из формулы (8) при $\tau\to 0+$ и $w\in E$ также является 


равномерно бесконечно малой величиной.


\end{pf}








Вариации вида (8) были получены ранее


в [6],\,\footnote{В статье [7] отсутствует ссылка на названную работу [6],


поскольку к моменту выхода статьи [7] автор данной


работы не был знаком с работой [6].


Автор благодарит профессора В.В.\,Старкова за указание на эту работу.}


но в более жестких 


предположениях относительно функции $Q(\zeta)$.


\medskip





{\bf 4.}


Исследование экстремальной ситуации с помощью вариаций (8)


для большого числа задач приводит к вариационному


условию, которое содержит








\begin{lem}


Пусть $\Gamma$ --- некоторый 


невырожденный континуум, принадлежащий 


$\xi$-плос\-ко\-сти, дополнение к которому состоит


из односвязной области, и для любой функции $p(\xi)$,


регулярной в некоторой проколотой окрестности


$\Dot{U}(\Gamma)$ континуума $\Gamma$ и 


имеющей вид $(7)$, выполняется


неравенство\,\footnote{Условие (18) возникает в экстремальных задачах, 


в которых функционалы выражаются через контурные 


интегралы от аналитических функций.}


\begin{equation}


\re\int_{\gamma}S(\xi)p(\xi)d\xi\le 0,


\end{equation}


где $\gamma$ -- произвольная замкнутая 


спрямляемая жорданова кривая,


лежащая в окрестности $\Dot{U}(\Gamma)$


и разделяющая граничные компоненты области


$\Dot{U}(\Gamma);$ интегрирование вдоль кривой


$\gamma$ осуществляется в направлении, при 


котором континуум $\Gamma$ остается слева$;$


$S(\xi)$ --- функция, регулярная в окрестности 


$\Dot{U}(\Gamma)$ континуума $\Gamma;$


$\xi(\zeta)$ --- функция, обратная для  


функции $\zeta(\xi)$ из теоремы $2$.





Тогда функция 


\begin{equation}


D(\zeta)=S(\xi(\zeta))(i\zeta \xi'(\zeta))^{2},


\end{equation}


регулярная в кольце $D(\delta,1)$, аналитически 


продолжается в кольцо $D(\delta^{2},1)$ и на


окружности $|\zeta|=\delta$ принимает вещественные


неотрицательные значения.


\end{lem}


\setcounter{footnote}{0}





\begin{pf}


Функция $\xi=\xi(\zeta)$ выполняет


конформное отображение кольца


$D(\delta,1)$ на область $\Dot{U}(\Gamma)$ так,


что граничной окружности $|\zeta|=\delta$ этого кольца


отвечает граничная компонента области $\Dot{U}(\Gamma)$,


принадлежащая континууму $\Gamma$.


В интеграле (18) с помощью подстановки


$\xi=\xi(\zeta)$ перейдем к переменной $\zeta$.


В качестве контура интегрирования


$\gamma$ выберем образ окружности 


$|\zeta|=\rho$, $\delta<\rho<1$,


при отображении $\xi=\xi(\zeta)$.


Учитывая (7), условие (18) можно записать в виде


\begin{equation}


\re \int\limits_{|\zeta|=\rho}D(\zeta)Q(\zeta)d\theta\ge 0,


\end{equation}


где $\theta=\arg\zeta$. В кольце $D(\delta,1)$ функцию $D(\zeta)$ 


можно разложить в ряд Лорана


\begin{equation}


D(\zeta)=\sum\limits_{-\infty}^{+\infty}c_{k}\zeta^{k}.


\end{equation}


Выберем функцию


\begin{equation}


Q(\zeta)=


\mathrm e^{i\psi}\bigg(\zeta^{n}-\frac{\delta^{2n}


\mathrm e^{-2i\psi}}{\zeta^{n}}\bigg),


\quad n=1,2,\dotsc,


\end{equation}


где $\psi$ --- произвольный вещественный параметр.


Реальная часть функции (22) обращается в тождественный


нуль на окружности $|\zeta|=\delta$.


Подставляя (21) и (22) в (20) и вычисляя


интеграл, получим


$$


\re (c_{-n}\mathrm e^{i\psi}-c_{n}\delta^{2n}\mathrm e^{-i\psi}~)\ge 0,


\quad n=1,2,\dotsc,


$$


причем это неравенство выполняется при любом значении


$\psi$. Отсюда


$$


c_{-n}=\delta^{2n}\overline{c}_{n},\quad n=1,2,\dotsc


$$


Следовательно, 


$$


D(\zeta)=c_{0}+\sum\limits_{k=1}^{+\infty}(c_{k}\zeta^{k}+


\delta^{2k}\overline{c}_{k}\zeta^{-k}).


$$


Из полученной формы разложения для $D(\zeta)$ 


видно, что ряд Лорана для функции $D(\zeta)$


сходится в кольце $D(\delta^{2},1)$


и $\im D(\zeta)\equiv \im c_{0}$ на окружности $|\zeta|=\delta$.





Легко показать, что $\im c_{0}=0$. Для этого в  


неравенстве (20) достаточно положить $Q(\zeta)=\mathrm e^{i\psi}$,


$\psi$ --- вещественный параметр, не зависящий от $\zeta$,


с областью изменения $[-\pi/2, +\pi/2]$, и проварьировать параметр 


$\psi$ в границах его изменения. 


Таким образом, функция $D(\zeta)$ регулярна в кольце


$D(\delta^{2},1)$, а на окружности $|\zeta|=\delta$


имеет мнимую часть, равную нулю.





Oстается показать, что на окружности


$|\zeta|=\delta$ вещественная часть функции


$D(\zeta)$ неотрицательна. С этой целью


в качестве функции $Q(\zeta)$ рассмотрим ядро Шварца


\begin{equation}


Q(\zeta)=\frac{\zeta+z}{\zeta-z},


\end{equation}


где $z$ --- произвольная точка открытого круга радиуса


$\delta$ с центром в начале, а $\zeta$ --- переменная точка


кольца $D(\delta,1)$. Легко проверяется, что функция


(23) регулярна в кольце $D[\delta,1)$ и имеет 


положительную вещественную часть на окружности $|\zeta|=\delta$,


поэтому на основании (20) интеграл Шварца


\begin{equation}


f(z)=\frac{1}{2\pi}


\int\limits_{|\zeta|=\rho}D(\zeta)\frac{\zeta+z}{\zeta-z}d\theta,


\quad \delta\le\rho<1,


\end{equation}


представляющий собой функцию, регулярную в круге $|z|\le\delta$, 


имеет в этом круге неотрицательную вещественную часть


$$


\re f(z)\ge 0,\quad |z|\le\delta.


$$


Так как на окружности $|z|=\delta$


мнимая часть функции $D(z)$ обращается тождественно в нуль,


из формулы (24) следует равенство


$$


D(z)=\re f(z),


$$


выполняющееся во всех точках окружности $|z|=\delta$.


Отсюда на окружности $|z|=\delta$ получаем неравенство


$D(z)\ge 0$.


\end{pf}





\begin{cornonum}


Нули функции


$D(\zeta)$, расположенные на окружности $|\zeta|=\delta$,


могут иметь только четный порядок, причем самих нулей может 


быть лишь конечное число.


\end{cornonum}





{\bf 5.}


Следующее утверждение двойственно основной лемме 


М.\,Шиффера [4] о граничных вариациях. 





\begin{lem}


В условиях леммы $1$, когда функция 


$S(\xi)\not\equiv 0$ регулярна в полной окрестности континуума 


$\Gamma$, граница дополнения к континууму $\Gamma$ до всей


плоскости представляет собой кусочно-аналитическую кривую, 


составленную из конечного числа траекторий квадратичного 


дифференциала $S(\xi)(d\xi)^{2}$ или их дуг, с точками 


ветвления лишь в нулях функции $S(\xi)$.


\end{lem}





\begin{pf}


Вначале покажем, что каждый простой конец области


$\Dot{U}(\Gamma)=U(\Gamma)\setminus\Gamma$, принадлежащий 


континууму $\Gamma$, состоит ровно из одной точки. 


Предположим противное. Тогда существует 


по меньшей мере один простой конец $P\subset\Gamma$, 


представляющий собой невырожденный континуум и отвечающий 


некоторой точке $\zeta'$ граничной окружности $|\zeta|=\delta$ 


кольца $D(\delta,1)$ при отображении $\xi=\xi(\zeta)$. 


Поскольку $S(\xi)\not\equiv 0$, с самого начала можно 


считать, что в полной


окрестности $U(\Gamma)$ континуума $\Gamma$ располагается 


лишь конечное число нулей этой


функции.\,\footnote{Иначе окрестность можно уменьшить.}


Каждую точку, изображающую нуль функции


$S(\xi)$, окружим столь малым открытым кружком с центром в 


соответствующей точке, чтобы кружки располагались внешним 


образом по отношению друг к другу и чтобы в дополнении 


к замыканию их объединения имелась непустая часть $P'$ 


континуума $P$. Обозначим через $T$ объединение названных кружков.


Пусть $p$ и $p'$ ---  различные точки множества $P'$. 


Выберем две последовательности $\{\xi_{n}\}$ и $\{\xi'_{n}\}$ 


точек области $\Dot{U}(\Gamma)$, сходящиеся соответственно к 


точкам $p$ и $p'$. В кольце $D(\delta,1)$ им соответствуют 


две последовательности $\{\zeta_{n}=\zeta(\xi_{n})\}$ и


$\{\zeta'_{n}=\zeta(\xi'_{n})\}$, сходящиеся к точке


$\zeta'$. Удалив из последовательностей $\{\zeta_{n}\}$


и $\{\zeta'_{n}\}$ конечное число членов, можно добиться, чтобы


среди элементов этих последовательностей не было точек


множества $F=\zeta(\Dot{U}(\Gamma)


\bigcap T^{*})$.\,\footnote{$T^{*}$ обозначает замыкание множества $T$.}





Далее будет описан алгоритм построения некоторой 


последовательности гладких дуг, стягивающихся к точке $\zeta'$.





Для каждого $n$ соединим точки $\zeta_{n}$ и $\zeta'_{n}$


дугой $\lambda_{n}$ окружности с центральным углом, 


не превосходящим $\pi$, целиком расположенной


в замкнутом круговом кольце $D[r_{n},R_n]$, где


$r_{n}=\min(|\zeta_{n}|, |\zeta'_{n}|)$ и 


$R_{n}=\max(|\zeta_{n}|, |\zeta'_{n}|)$. 


При $n\to\infty$ дуги $\lambda_{n}$ равномерно стягиваются к точке $\zeta'$.


Если в последовательности $\{\lambda_{n}\}$ существует


подпоследовательность $\{\lambda_{n_{k}}\}$ дуг, не


пересекающихся с множеством $F$, выбираем эту


подпоследовательность и на этом останавливаемся,


в противном случае строим новую последовательность


дуг следующим образом. Возьмем дугу $\lambda_{n}$,


имеющую непустое пересечение с множеством $F$, и 


будем двигаться по ней от точки $\zeta_{n}$ до тех пор, 


пока впервые не встретим точку $\zeta''_{n}\in F$.


После этого заменим дугу $\lambda_{n}$ ее частью $\lambda'_{n}$,


расположенной между точками $\zeta_{n}$ и $\zeta''_{n}$.


В результате получим новую последовательность дуг


$\{\lambda'_{n}\}$, стягивающуюся при $n\to\infty$ к


точке $\zeta'$. При $n\to\infty$ концы $\zeta''_{n}$


построенных дуг стремятся к $\zeta'$, поэтому все


предельные точки последовательности 


$\{\xi''_{n}=\xi(\zeta''_{n})\}$ принадлежат


множеству $P\bigcap T^{*}$. Из последовательности 


$\{\xi''_{n}\}$ извлекаем подпоследовательность 


$\{\xi''_{n_{k}}\}$, сходящуюся к некоторой точке


$p''\in P\bigcap T^{*}$. Очевидно, $p''\not= p$, поскольку


$p\notin T^{*}$. Теперь из последовательности 


$\{\lambda'_{n}\}$ отбираем дуги с концами в точках


$\zeta''_{n_{k}}=\zeta(\xi''_{n_{k}})$. В результате


получаем последовательность дуг $\{\lambda'_{n_{k}}\}$.


На этом алгоритм построения дуг заканчивается. 





В результате осуществления описанного алгоритма


в любом случае приходим к некоторой последовательности


дуг $\{\lambda_{k}\}$, раположенных в кольце $D(\delta,1)$,


равномерно стягивающихся к точке $\zeta'$


и обладающих следующими свойствами: 


\begin{itemize}


\item[a)] любая дуга $\lambda_{k}$ 


представляет собой дугу окружности с центральным углом, 


не превышающим $\pi$; 


\item[b)] дуги $l_{k}=\xi(\lambda_{k})$ лежат в множестве


$\Dot{U}(\Gamma)\setminus T$, а их концы $\xi_{1,k}$ и 


$\xi_{2,k}$ стремятся к различным точкам $p_{1}$ и 


$p_{2}$ простого конца $P$. 


\end{itemize}


\setcounter{footnote}{0}





\noindent{}Обозначим через $d$ расстояние между точками 


$p_{1}$ и $p_{2}$. Очевидно, $d>0$.  


При $k\to\infty$ дуги $l_{k}$ равномерно приближаются к 


$P$. Положим


$L=\bigcup\limits_{k}l_{k}$.


Замыкание $L^{*}$ множества $L$ отделено от нулей функции 


$S(\xi)$, поэтому модуль функции $S(\xi)$ на множестве $L^{*}$ 


имеет положительную нижнюю грань $m>0$. Между тем, 


прообразы $\lambda_{k}$ дуг $l_{k}$ при 


отображении $\xi=\xi(\zeta)$, начиная с некоторого 


номера, принадлежат произвольно малой $\varepsilon$-окрестности 


точки $\zeta'$ и, следовательно, имеют диаметр, не превосходящий 


$2\varepsilon$. Наконец, в соответствии с леммой 1 модуль функции 


$D(\zeta)$ в кольце $D[\delta,\delta+\varepsilon)$ 


ограничен сверху некоторой положительной постоянной $M$. 


В силу сказанного, а также ввиду соотношения (19), 


начиная с некоторого номера, на дугах $\lambda_{k}$ 


должно выполняться неравенство


$$


|\xi'(\zeta)|\le \bigg(\frac{M}{\delta^{2}m}\bigg)^{\frac{1}{2}}.


$$


С другой стороны,


$$


\xi_{2,k}-\xi_{1,k}=\int_{\lambda_{k}}\xi'(\zeta)d\zeta,


$$


поэтому, начиная с некоторого номера, величина 


$|\xi_{2,k}-\xi_{1,k}|$ не превосходит 


$2\pi\varepsilon(M/\delta^{2} m)^{1/2}$. При достаточно малом


$\varepsilon$ это меньше, чем $d/2$, что невозможно, 


т.\,к. точки $\xi_{1,k}$ и $\xi_{2,k}$ при $k\to\infty$


стремятся к $p_{1}$ и $p_{2}$ соответственно.





Итак, каждый простой конец границы области $\Dot{U}(\Gamma)$, 


принадлежащий континууму $\Gamma$, состоит из одной граничной 


точки. Следовательно, граничная компонента $\Omega$ окрестности 


$\Dot{U}(\Gamma)$, принадлежащая континууму $\Gamma$, состоит 


лишь из достижимых граничных точек. Это означает, что функцию


$\xi(\zeta)$ можно по непрерывности продолжить на


граничную окружность $|\zeta|=\delta$ кольца 


$D(\delta,1)$, так что в итоге функция $\xi(\zeta)$ 


будет непрерывной в кольце $D[\delta,1)$,


а соответствие достижимых точек области $\Dot{U}(\Gamma)$,


составляющих множество $\Omega$, и точек окружности 


$|\zeta|=\delta$ будет взаимно однозначным.   


Прообразы лежащих на $\Omega$ нулей функции $S(\xi)$ 


при отображении $\xi=\xi(\zeta)$ на окружности $|\zeta|=\delta$ 


образуют замкнутое множество, дополнение к которому состоит


из не более чем счетного числа открытых дуг. 


Выберем любую из них. Пусть это будет


дуга $\theta'<\theta<\theta''$. 





Покажем, что на выбранной


дуге функция $\xi(\zeta)$ имеет непрерывную производную по


переменной $\theta=\arg\zeta$. С этой целью внутри интервала


$(\theta',\theta'')$ выберем произвольный отрезок 


$[\theta_{0},\theta_{1}]$. Ограничимся значениями 


$\rho$, $\delta<\rho<\rho^{*}$, столь близкими к $\delta$,


что на окружности $|\zeta|=\rho$ нет нулей функции


$S(\xi(\zeta))$. На дуге $\theta_{0}\le\theta\le\theta_{1}$


окружности $|\zeta|=\rho$ для соответствующим образом подобранной


ветви квадратного корня равенство (19) можно переписать в виде


\begin{equation}


\frac{d\xi}{d\theta}=\bigg(\frac{D(\zeta)}{S(\xi)}\bigg)^{\frac{1}{2}},


\end{equation}


откуда интегрованием получим равенство


$$


\xi(\zeta)-\xi(\zeta_{0})=\int_{\theta_{0}}^{\theta}


\bigg (\frac{D(\zeta)}{S(\xi)}\bigg )^{\frac{1}{2}}d\arg\zeta,


$$


выполняющееся на дуге $\theta_{0}\le\theta\le\theta_{1}$


окружности $|\zeta|=\rho$.


Путем предельного перехода при $\rho\to\delta+0$ убеждаемся,


что последнее равенство справедливо и на дуге 


$\theta_{0}\le\theta\le\theta_{1}$ окружности $|\zeta|=\delta$. 


Поскольку подинтегральное выражение непрерывно зависит от $\arg\zeta$,


функция $\xi(\zeta)$ действительно непрерывно дифференцируема


по $\theta=\arg\zeta$ внутри дуги $[\theta_{0},\theta_{1}]$ окружности


$|\zeta|=\delta$. Так как промежуток $[\theta_{0},\theta_{1}]$ был 


произвольно выбран внутри интервала $(\theta',\theta'')$, 


то действительно функция $\xi(\zeta)$ имеет непрерывную


производную по $\theta=\arg\zeta$ на дуге $\theta'<\theta<\theta''$


окружности $|\zeta|=\delta$. 





Итак, функция $\xi(\zeta)$


непрерывно дифференцируема внутри любой дуги окружности 


$|\zeta|=\delta$ с концами в соседних прообразах нулей


функции $S(\xi)$. Это означает, что любой участок 


такой дуги, на котором нет нулей $D(\zeta)$, функция 


$\xi=\xi(\zeta)$ отображает в кривую класса $C^{1}$. 


Заметим, что равенство (19) 


можно переписать в виде


$S(\xi)(d\xi)^{2}=D(\zeta)(d\theta)^{2}$,


где $|\zeta|=\delta$ и $\theta=\arg\zeta$. 


Следовательно, гладкие дуги, входящие в состав 


$\Omega$, служат частями траекторий квадратичного 


дифференциала $S(\xi)(d\xi)^{2}$,  


поэтому функция $\xi(\zeta)$ является аналитической 


в любой внутренней точке $\zeta_{0}$ произвольной 


дуги окружности $|\zeta|=\delta$ 


с концами в соседних прообразах нулей функции $S(\xi)$. 


Действительно, в силу леммы 1 функция


$$


\varkappa(\zeta)=


-i \int_{\zeta_{0}}^{\zeta}D^{\frac{1}{2}}(\zeta')


\frac{d\zeta'}{\zeta'}


$$


регулярна в точке $\zeta_{0}$.


Утверждение об аналитичности $\xi(\zeta)$ в 


окрестности точки $\zeta_{0}$ вытекает из  представления


$\xi(\zeta)=\mu^{-1}\circ\varkappa(\zeta)$,\,\footnote{Это представление 


является следствием равенства (25).}


где  


$$


\mu(\xi)=\int_{\xi_{0}}^{\xi}S^{\frac{1}{2}}(\xi)d\xi,\quad


\xi_{0}=\xi(\zeta_{0}),


$$


\setcounter{footnote}{0}





\noindent{}--- функция, регулярная и однолистная 


в малой окрестности точки $\xi_{0}$,


а $\mu^{-1}$ --- функция, обратная к $\mu(\xi)$.





Если $D(\zeta_{0})\ne0$, то в силу (25) \ $\xi(\zeta)$ 


однолистна в некоторой окрестности точки $\zeta_{0}$, и 


монотонному движению точки $\zeta$ вдоль окружности 


$|\zeta|=\delta$, происходящему в пределах малой 


окрестности точки $\zeta_{0}$, отвечает монотонное 


движение точки $\xi=\xi(\zeta)$ вдоль траектории


квадратичного дифференциала $S(\xi)(d\xi)^{2}$. 


Если $D(\zeta_{0})=0$, то в некоторой окрестности


точки $\zeta_{0}$ функция $D(\zeta)$ не будет иметь других


нулей. Ввиду следствия из леммы 1, равенства (25), 


а также однолистности функции 


$\xi(\zeta)$ в кольце $D(\delta,1)$ точка $\zeta_{0}$ 


может быть лишь нулем второго порядка для $D(\zeta)$. 


Монотонному движению точки $\zeta$ по дуге окружности 


$|\zeta|=\delta$, происходящему {\it до} и {\it после} 


точки $\zeta_{0}$ в пределах малой окрестности точки $\zeta_{0}$, 


в этом случае также отвечает монотонное движение точки $\xi=\xi(\zeta)$ 


вдоль траектории квадратичного дифференциала $S(\xi)(d\xi)^{2}$, но при


переходе точки $\zeta$ через $\zeta_{0}$ направление движения


точки $\xi=\xi(\zeta)$ вдоль траектории изменяется на 


противоположное.


Следовательно, нули функции $D(\zeta)$, расположенные на окружности


$|\zeta|=\delta$ и отличные от прообразов нулей $S(\xi)$, 


имеют второй порядок и соответствуют концевым точкам аналитических


дуг континуума $\Omega$.  





Вновь рассмотрим произвольную дугу $(\theta', \theta'')$


окружности $|\zeta|=\delta$ с концами в соседних прообразах


нулей функции $S(\xi)$. Монотонному движению точки


$\zeta$ по дуге $(\theta', \theta'')$ будет соответствовать


монотонное движение ее образа $\xi=\xi(\zeta)$ по траектории


квадратичного дифференциала $S(\xi)(d\xi)^{2}$, если только


при своем движении точка $\zeta$ не будет встречать нулей


функции $D(\zeta)$. При прохождении точки $\zeta$ через 


нуль функции $D(\zeta)$, расположенный на дуге


$(\theta', \theta'')$, движение точки $\xi$ будет продолжаться по 


той же траектории, но в противоположном направлении.


По этой причине между соседними прообразами нулей функции


$S(\xi)$ может располагаться не более одного нуля


функции $D(\zeta)$.\,\footnote{Может случиться, что на $\Omega$ совсем нет


нулей функции $S(\xi)$, и тогда либо на окружности 


$|\zeta|=\delta$ вовсе нет нулей $D(\zeta)$, и кривая


$\Omega$ представляет собой замкнутую траекторию


дифференциала $S(\xi)(d\xi)^{2}$, либо на окружности


$|\zeta|=\delta$ имеется ровно два нуля второго


порядка функции $D(\zeta)$, и тогда кривая $\Omega$


является разомкнутой дугой траектории дифференциала 


$S(\xi)(d\xi)^{2}$. Впрочем, если $\Omega$ представляет 


собой замкнутую траекторию, то $S(\xi)\equiv 0$.} 


%%%%%


Следовательно, дуга 


$(\theta', \theta'')$ окружности $|\zeta|=\delta$


при отображении $\xi=\xi(\zeta)$ переводится


либо в траекторию дифференциала $S(\xi)(d\xi)^{2}$,


имеющую концы в нулях функции $S(\xi)$, либо в


дугу траектории дифференциала $S(\xi)(d\xi)^{2}$,


исходящую из нуля функции $S(\xi)$. Таких дуг


может существовать только конечное число. Каждая


такая дуга соответствует одной открытой дуге окружности


$|\zeta|=\delta$ с концами в соседних прообразах


нулей функции $S(\xi)$. Поэтому


дополнение ко множеству прообразов нулей функции $S(\xi)$


на окружности $|\zeta|=\delta$ состоит из конечного числа


открытых дуг. Так как $\xi(\zeta)\not\equiv \mathrm{const}$,


число прообразов нулей функции $S(\xi)$, 


расположенных на окружности $|\zeta|=\delta$, также


конечно. В противном случае на окружности $|\zeta|=\delta$


существовала бы целая дуга, целиком отображающаяся


в какой-либо нуль функции $S(\xi)$.


\end{pf}





В заключение хочу высказать глубокую признательность профессору 


Л.А.\,Аксентьеву за доверие и поддержку, а также поблагодарить рецензента, 


указавшего на ошибку в первоначальной редакции доказательства


леммы 2.








\begin{thebibliography}{99}





\bibitem{1}


Ерохин В.Д. {\em О конформных преобразованиях колец и об 


основном базисе пространства функций, аналитических в 


элементарной окрестности континуума} //


ДАН СССР. -- 1958. -- Т.\,120. -- \No\,4. -- С.\,689--696.


\bibitem{2}


Гельфер С.А. {\em О распространении вариационного метода


Голузина--Шиффера на многосвязные области} //


ДАН СССР. -- 1962. -- Т.\,142. -- \No\,3. -- С.\,503--506.


\bibitem{3}


Голузин Г.М. {\em Метод вариаций в конформном отображении.} I //


Матем. сб. -- 1946. -- Т.\,19. -- \No\,2. -- С.\,203--236.


\bibitem{4}


Schiffer M. {\em A method of variation within the family of


simple functions} // Proc. London Math. Soc. -- 1938. 


-- V.\,44. -- P.\,432--449.


\bibitem{5}


Schiffer M. {\em Variation of the Green function and the theory


of p-valued functions} // Amer. J. Math. -- 1943. -- V.\,65. --


P.\,341--360.


\bibitem{6}


Бабенко К.И. {\em К теории экстремальных задач для 


однолистных функций класса $S$} // Тр. Матем. ин-та АН СССР. --


М.: Наука. -- 1972. -- Т.\,101. -- 320 с.


\bibitem{7}


Кожевников В.В. {\em Факторизация конформных отображений


и вариации однолистных функций} // Докл. РАН. -- 2000. -- Т.\,375. -- 


\No\,3. -- С.\,305--306.





\end{thebibliography}





\vskip 1cm





\postup{\it Кубанский государственный}{\qquad \it Поступили}


\postup{\it университет}{\it первый вариант $21.03.2002$}


\postup{}{\it окончательный вариант $31.03.2003$}





\label{end}


\end{document}


